\documentclass[10pt]{article}
\usepackage[T1]{fontenc}
\usepackage[utf8]{inputenc}
\usepackage[french]{babel}
\usepackage{amsmath,amssymb,mathrsfs,mathpazo}
\usepackage{geometry}
\geometry{a4paper,margin=2cm}
\usepackage{graphicx,array,tabularx,tabularray,colortbl,multirow,multicol,booktabs,enumitem,pifont,ulem,fancybox,diagbox,cancel}
\usepackage{pstricks,pstricks-add,pst-plot,pst-tree,pst-node,pst-text,pst-func,pst-math,pst-bspline,pst-3dplot,pst-eucl,pst-solides3d}
\let\pscalculate\undefined
\usepackage{pst-bezier}
\usepackage{tikz,pgfplots,tkz-tab}
\usetikzlibrary{arrows,arrows.meta,patterns,calc,positioning,trees,decorations.markings,shapes.geometric,shapes.misc,shapes.arrows,intersections,decorations.pathmorphing,decorations.pathreplacing,backgrounds,shadows,babel}
\pgfplotsset{compat=1.18}
\usepgfplotslibrary{fillbetween}
\usepackage[np]{numprint}
\usepackage{listings,makecell,pas-tableur,twoopt}
\let\barre\undefined
\usepackage[pstricks]{tablvar}
\definecolor{darkgreen}{RGB}{0,100,0}
\providecommand{\red}{\color{red}}
\providecommand{\blue}{\color{blue}}
\providecommand{\green}{\color{green}}
\providecommand{\black}{\color{black}}
\providecommand{\white}{\color{white}}
\providecommand{\gray}{\color{gray}}
\providecommand{\cyan}{\color{cyan}}
\providecommand{\magenta}{\color{magenta}}
\providecommand{\yellow}{\color{yellow}}
\providecommand{\violet}{\color{violet}}
\providecommand{\orange}{\color{orange}}
\providecommand{\euro}{\texteuro}
\providecommand{\EUR}{\texteuro}
\providecommand{\ieme}{\textsuperscript{e}}
\providecommand{\degres}{\ensuremath{^\circ}}
\providecommand{\degree}{\ensuremath{^\circ}}
\providecommand{\vv}[1]{\overrightarrow{#1}}
\providecommand{\Coord}[1]{#1}
\providecommand{\decosix}{}
\providecommand{\starredbullet}{\ensuremath{\bullet}}
\pagestyle{empty}
\setlength{\parindent}{0pt}
\usepackage{adjustbox}\DeclareUnicodeCharacter{2236}{:}\DeclareUnicodeCharacter{2190}{\ensuremath{\leftarrow}}
\providecommand{\R}{\mathbb{R}}
\providecommand{\N}{\mathbb{N}}
\providecommand{\Z}{\mathbb{Z}}
\providecommand{\Q}{\mathbb{Q}}
\providecommand{\D}{\mathbb{D}}
\providecommand{\C}{\mathbb{C}}
\providecommand{\e}{{\mathrm{e}}}
\providecommand{\vect}[1]{\overrightarrow{#1}}
\providecommand{\vectt}[1]{\overrightarrow{\mathrm{#1}}}
\providecommand{\vvt}[1]{\overrightarrow{\mathrm{#1}}}
\providecommand{\np}[2][]{{#2~#1}}
\providecommand{\vv}[1]{{\overrightarrow{#1}}}
\renewcommand{\d}{{\mathrm{d}}}
\providecommand{\degres}{{^\circ}}
\providecommand{\degre}{{^\circ}}
\providecommand{\blue}{\color{blue}}
\providecommand{\red}{\color{red}}
\providecommand{\enskip}{\quad}
\providecommand{\llbracket}{\lbrack\!\lbrack}
\providecommand{\rrbracket}{\rbrack\!\rbrack}
\providecommand{\euro}{{\text{€}}}
\providecommand{\no}{{\text{n°}}}
\providecommand{\interval}[2]{{[#1;#2]}}
\providecommand{\intervFO}[2]{{[#1;#2[}}
\providecommand{\dotfill}{\ldots}

\providecommand{\ds}{\displaystyle}
\providecommand{\textding}[1]{#1}
\providecommand{\Oij}{\(\mathrm{O};\vec{\imath},\vec{\jmath}\)}
\providecommand{\Oijk}{\(\mathrm{O};\vec{\imath},\vec{\jmath},\vec{k}\)}
\providecommand{\Ouv}{\(\mathrm{O};\vec{u},\vec{v}\)}
\providecommand{\intervFF}[2]{\ensuremath{[#1;#2]}}
\renewcommand{\euro}{\ifmmode\text{\texteuro}\else\texteuro\fi}
\DeclareUnicodeCharacter{00B0}{\ifmmode^{\circ}\else\textdegree\fi}
\DeclareUnicodeCharacter{20AC}{\ifmmode\text{\texteuro}\else\texteuro\fi}
% Énoncés : Éducation nationale. Transcription et corrigés : APMEP, Denis Vergès et contributeurs. Découpage et classement : CoopMaths, Sébastien Lozano, Éric Elter et contributeurs. Adaptation de la présentation : Maths Explorer.
% CC BY-SA : https://coopmaths.fr/www/about/
% https://coopmaths.fr/alea/static/eam/2026/tex/eam_specialite_2026_06_antilles_1.tex
\begin{document}
\noindent\textit{APMEP / CoopMaths. CC BY-SA. Présentation adaptée par Maths Explorer.}\par\medskip
On s'intéresse à l'évolution de la valeur d'une voiture.

En janvier 2025, la valeur de la voiture est égale à 10\,000 euros.

On suppose qu'ensuite, chaque année, la valeur diminue de 10\,\%.

Pour tout entier naturel $n$, on note $u_n$ la valeur de la voiture, en euros, au mois de janvier de l'année $2025 + n$.

On a donc $u_0 = \np{10000}$.

\begin{enumerate}
	\item \begin{enumerate}
		\item Calculer la valeur de la voiture en janvier 2026.
		\item Justifier que pour tout entier naturel $n$, on a
		\[u_{n+1} = 0,9\, u_n.\]
	\end{enumerate}

	\item En déduire la nature de la suite $(u_n)$ et préciser sa raison.

	\item \begin{minipage}[t]{0.75\linewidth}
		\begin{enumerate}
			\item Pour tout entier naturel $n$, exprimer le terme $u_n$ en fonction de $n$.

			\item En déduire la valeur de la voiture en janvier 2030.

			\emph{On pourra s'aider de l'aide au calcul ci-contre.}
		\end{enumerate}
	\end{minipage}
	\hfill
	\begin{tblr}[t]{|X[3.5cm]|}\hline
		Aide au calcul\\
		$0,9^4 = 0,6561$\\
		$0,9^5 = 0,59049$\\
		$0,9^6 = 0,531441$\\ \hline
	\end{tblr}

	\item \begin{minipage}[t]{0.62\linewidth}
		On considère le programme Python ci-contre.

		Recopier et compléter ce programme afin qu'il renvoie, après exécution, la plus petite valeur de $n$ telle que \quad$u_n < A$,\quad où $A$ est un nombre réel strictement positif.
	\end{minipage}\hfill
	\begin{minipage}[t]{5.5cm}
		\begin{lstlisting}[style=python,numbers=none,xleftmargin=0mm,xrightmargin=0mm]
def seuil(A) :
	N = 0
	U = 10000
	while ......... :
		U = .......
		N = N+1
	return N
\end{lstlisting}
	\end{minipage}

	\item L'exécution de ce programme, pour plusieurs valeurs de $A$, donne les résultats suivants :

\begin{lstlisting}[style=python,xleftmargin=5cm,xrightmargin=5cm,numbers=none]
>>> seuil(7500)
  3
>>> seuil(5000)
  7
>>> seuil(2500)
  14
>>> seuil(2000)
  16
\end{lstlisting}

	\begin{enumerate}
		\item Interpréter, dans le contexte de l'exercice, le résultat obtenu lors de l'appel \texttt{seuil(5000)}.

		\item Déterminer, en expliquant la démarche, l'année à partir de laquelle la voiture aura perdu plus des trois quarts de sa valeur initiale.
	\end{enumerate}
\end{enumerate}

\end{document}
